\scanpage{166}\mspage{111}
% archive sheet number: 165
essentiellement le foncteur déduit de l'inclusion
\[
\left(\underline{\mathrm{Fib}}\,S,\underbrace{W_S\mid\underline{\mathrm{Fib}}\,S}_{\shortparallel\,W_S^{\unsure{0|o|\circ}}\mid\underline{\mathrm{Fib}}\,S}\right)\longrightarrow(\unsure{M|\mathcal M}_{/S},W_S^0)
\]
par localisation au \uline{moyen de} $W_S^0$, donc le \deleted{\illegible{1}} foncteur
\[
(\underline{\mathrm{Fib}}\,S)(W_S^0)^{-1}\longrightarrow(M_{/S})(W_S^0)^{-1},\tag*{$(**)$}
\]
compte tenu de l'équivalence $(*)$ de \textcircled{I} ; i.e. $i$ est le composé du quasi-inverse de $(*)$, avec le foncteur $(**)$ précédent. Explicitement, il est défini \unsure{sur|par}
\[
i(X)=\text{Classe de }C(X)\text{ dans }(M_{/S})(W_S^0)^{-1}\,;
\]
en tenant compte du fait que $X\longmapsto C(X)$ transforme $W_S$ en $W_S^0$, grâce à la condition 3°) sur $W_S^0$.

Ceci posé, \uline{sous l'hyp. B} (``Quillen''), on a : dans le couple de foncteurs $(\rho,i)$ de $(*)$ page 110, \uline{$i$ est adjoint à droite de $\rho$}, i.e. on a pour $\xi\in\mathrm{Ob}\,\mathrm{Ho}_{W_S^0}(S)$, $\eta\in\mathrm{Ho}^{\mathrm{lc}}_W(S)$
\[
\mathrm{Hom}_{W_S^0}(\xi,i(\eta))\simeq\mathrm{Hom}_{W_S}(\rho(\xi),\eta),
\]
ce qui s'explicite, pour $X,Y$ dans $\deleted{\unsure{\in}}M_{/S}$, par
\[
\boxed{\mathrm{Hom}_{W_S^0}(X,C(Y))\simeq\mathrm{Hom}_{W_S}(X,Y)}
\]
% slanted
\margin{\uline{NB} le premier membre est indépendant du choix de $W_S^0$ ! (Pourvu seulement que l'hyp $W_S^0\cap\mathrm{Fl}\,\underline{\mathrm{Fib}}\,S=W_S\cap\mathrm{Fl}\,\underline{\mathrm{Fib}}\,S$, en plus de $W_S^0\subset W_S$, soit satisfaite.)}
\unsure{S'ensuit|Il s'ensuit} que \uline{le foncteur $i$ est pleinement fidèle.}

La même formule d'adjonction, et la
